% from ega4-21c-fr.tex
\oldpage[IV-4]{322}
La condition $b)$ peut encore s'exprimer de la façon suivante :

\emph{$b')$ Pour toute partie fermée $Y$ de $X$ telle que
$\operatorname{codim}(Y,X)\geq2$ et telle qu'il existe $\lambda$ pour lequel
$Y\subset X-U_\lambda$, le couple $(X,Y)$ est parafactoriel.}

La propriété $a)$ peut encore s'exprimer de la façon suivante : si l'ensemble fermé
$N$ des points $x\in X$ où $Z$ n'est pas principal est contenu dans l'un des
$X-U_\lambda$, alors $Z$ est localement principal. Comme $X$ est normal, la condition
$\dim(\mathcal O_{X,x})\leq1$ entraîne que $\mathcal O_{X,x}$ est un corps ou un
anneau de valuation discrète, donc $Z_x$ est principal; en d'autres termes, on a
nécessairement $\operatorname{codim}(N,X)\geq2$ (5.1.3). La propriété $a)$ est donc
équivalente à la suivante :

\emph{$a')$ S'il existe un ensemble fermé $Y$ contenu dans un des $X-U_\lambda$, tel
que $\operatorname{codim}(Y,X)\geq2$ et que $Z|(X-Y)$ soit localement principal,
alors $Z$ est localement principal.}

Notons d'autre part que $b)$ implique $b')$ en vertu de (21.13.10); inversement, si
$b')$ est vérifiée et si $x\notin U_\lambda$, alors $Y=\{x\}$ est contenue dans
$X-U_\lambda$ et on a $\operatorname{codim}(Y,X)\geq2$, donc il résulte encore de
(21.13.10) que $\mathcal O_{X,x}$ est parafactoriel, ce qui prouve l'équivalence de
$b)$ et $b')$. On est ainsi ramené à prouver l'équivalence de $a')$ et de $b')$.
Notons que puisque $X$ est normal, on a, pour toute partie fermée $Y$ de $X$ telle que
$\operatorname{codim}(Y,X)\geq2$, $\operatorname{prof}_Y(\mathcal O_X)\geq2$
(5.8.6); si l'on pose $U=X-Y$, l'homomorphisme
$\mathcal O_X\to j_*(\mathcal O_U)$ est donc bijectif (21.13.4); comme les conditions
$a')$ et $b')$ sont locales, on voit qu'il suffit de montrer l'équivalence des
conditions suivantes, lorsque $X$ est noethérien et normal et $Y$ fermé dans $X$ et
tel que $\operatorname{codim}(Y,X)\geq2$ :

\emph{$a'')$ Pour tout cycle $1$-codimensionnel $Z$ sur $X$, l'hypothèse que $Z|U$
(où $U=X-Y$) est localement principal entraîne que $Z$ lui-même est localement
principal.}

\emph{$b'')$ L'homomorphisme canonique
$\operatorname{Pic}(X)\to\operatorname{Pic}(U)$ est surjectif.}

Prouvons d'abord que $a'')$ entraîne $b'')$; un élément
$\zeta_0\in\operatorname{Pic}(U)$ a pour image dans $\mathfrak{Cl}(U)$ (21.6.11) la
classe d'un cycle $1$-codimensionnel $Z_0$ sur $U$, qui est localement principal.
L'hypothèse $\operatorname{codim}(Y,X)\geq2$ entraîne que l'homomorphisme de
restriction $\mathscr Z^1(X)\to\mathscr Z^1(U)$ est bijectif, donc $Z_0=Z|U$, où $Z$
est un cycle $1$-codimensionnel sur $X$; puisque $Z_0$ est localement principal, il en
est de même de $Z$ en vertu de $a'')$; il résulte de (21.6.11) que l'image de $Z$ dans
$\mathfrak{Cl}(X)$ est l'image d'un unique élément $\zeta\in\operatorname{Pic}(X)$,
et il est clair alors que $\zeta_0$ est l'image de $\zeta$.

Inversement, prouvons que $b'')$ entraîne $a'')$. Soit donc $Z$ un cycle
$1$-codimensionnel sur $X$ tel que $Z|U$ soit localement principal; l'image de $Z|U$
dans $\mathfrak{Cl}(U)$ est donc l'image d'un unique élément
$\zeta_0\in\operatorname{Pic}(U)$ (21.6.11). Par hypothèse il existe
$\zeta\in\operatorname{Pic}(X)$ dont l'image dans $\operatorname{Pic}(U)$ est
$\zeta_0$; l'image de $\zeta$ dans $\mathfrak{Cl}(X)$ est donc la classe d'un cycle
$1$-codimensionnel $Z'$ sur $X$ tel que $Z|U$ et $Z'|U$ soient linéairement
équivalents. Mais puisque $U$ est schématiquement dense dans $X$ l'image de
$\mathscr Z^1_{\mathrm{princ}}(X)$ par l'isomorphisme
$\mathscr Z^1(X)\to\mathscr Z^1(U)$ est $\mathscr Z^1_{\mathrm{princ}}(U)$, donc $Z$
et $Z'$ sont linéairement équivalents, et puisque $Z'$ est localement principal, il
en est de même de $Z$.
\end{proof}

\begin{corollary}[21.13.15]
\label{IV.21.13.15}
Soient $X$ un préschéma localement noethérien et normal, $S$ une partie de $X$. Les
conditions suivantes sont équivalentes :
\begin{enumerate}[label=\emph{\alph*)}]
  \item Tout cycle $1$-codimensionnel sur $X$ qui est principal aux points de $S$ est
  localement principal.
  